\documentclass[11pt]{book}

\usepackage{ppbook}

\begin{document}

\setcounter{chapter}{4}
\chapter{LLM for Specification Synthesis}

LLMs can make mistakes.

They also sometimes don't pay attention to the important
part of a long prompt.

As context gets very long, the middle part may get ignored.

Also, LLMs are autoregressive, so an error in one output
can become input for the next step:

\[
\text{error}_1
\rightarrow
\text{input}_2
\rightarrow
\text{error}_2
\rightarrow
\cdots
\]

So again the idea is:

\[
\text{LLM generates candidate spec}
\rightarrow
\text{static analysis / formal verification}
\rightarrow
\text{check}
\]

\section{AutoSpec}

Main idea = don't give the whole program to the LLM at once.

Break the program into smaller pieces first.

\[
\text{large program}
\rightarrow
\text{decomposition}
\rightarrow
\text{smaller parts}
\]

Program decomposition gives a call graph /
inter-procedural CFG.

Then specs are generated bottom-up.

If $f$ calls $f_1,f_2,\ldots$

first get specs for the callees,

then use those while generating the spec for $f$.

\[
f_1,f_2 \rightarrow f
\]

This reduces the amount of context the LLM has to
deal with at one time.

After generating a spec, don't directly trust it.

First check whether the spec is satisfiable / valid.

Then another check:

is this spec actually enough to finish the
verification?

If not,

\[
\text{strengthen spec}
\rightarrow
\text{verify again}
\]

so generation is iterative.

\section{Few-shot prompting}

Giving a few relevant examples helps.

Instead of only telling the LLM the task,
give examples like

\[
(\text{program},\text{corresponding spec})
\]

This helps it understand both the format and
how specifications should look.

For nested loops, legality of generated specs can
be checked early, while satisfiability may be
checked later.

Frama-C is used for verification in AutoSpec.

The main idea is still:

\[
\text{generate}
\rightarrow
\text{check}
\rightarrow
\text{strengthen if needed}
\]

\section{SpecGen}

Another approach for generating specifications.

The first phase is more conversation based.

Start with few-shot examples and ask the LLM to
generate specs.

Then the verifier gives feedback when something fails.

That failure information is used in the next
LLM round.

\[
Spec_0
\rightarrow
Verifier
\rightarrow
feedback
\rightarrow
LLM
\rightarrow
Spec_1
\]

and repeat.

\begin{ppnote}
Verifier errors can be too large or too abstract.

Instead of dumping everything to the LLM,
only one verification failure is reported per attempt.

Failures can then be collected and summarised into
some guidance.
\end{ppnote}

\section{Mutation}

Sometimes the LLM-generated spec is close,
but not exactly correct.

Instead of generating everything again from scratch,
take the failing specification and modify it.

\[
Spec
\rightarrow
\text{mutation}
\rightarrow
Spec'
\]

There are different mutation operators.

For example, logical operators can be changed,
comparisons can be changed, etc.

Then check which mutated specifications can
pass verification.

So:

\[
\text{failing spec}
\rightarrow
\text{mutate}
\rightarrow
\text{candidate specs}
\rightarrow
\text{verification}
\]

If a mutation still fails, it can be mutated again.

So this becomes a search over nearby specifications.

\section{Verification feedback}

Verification failure feedback is useful,
but error messages are often very abstract/general.

So the LLM may not understand exactly what needs
to be changed.

Mutation gives another way to explore possible
corrections.

For selecting mutations, candidates can be grouped
according to the same template / family.

Then a heuristic can be used to decide which
candidate should actually be sent for verification.

So not every possible mutation needs to be checked.

\section{Overall idea}

Both approaches use the same basic idea:

\begin{ppkey}
\[
\boxed{
\text{LLM generates}
\quad\rightarrow\quad
\text{formal tool checks}
}
\]

LLM is useful for producing possible specifications,
but verification is still done by the formal tools.
\end{ppkey}

AutoSpec:

\[
\text{decompose program}
\rightarrow
\text{generate bottom-up}
\rightarrow
\text{validate}
\rightarrow
\text{strengthen}
\]

SpecGen:

\[
\text{generate}
\rightarrow
\text{get failure feedback}
\rightarrow
\text{revise / mutate}
\rightarrow
\text{verify}
\]

The important thing is keeping the LLM in the
candidate-generation/search part rather than
letting its output alone decide correctness.

\end{document}
